
by putting, for every @@M0@@,
@@M1@@

\sgaStableTarget{sga6.e10.rem.5-2-2}\textbf{Remarks 5.2.2.}:
\begin{enumerate}[label=\alph*)]
\item Suppose that @@M2@@ is the locally ringed topos of Zariski sheaves on a scheme @@M3@@. When @@M4@@ is a Cartier divisor @@M5@@ on @@M6@@ and
@@M7@@
one has
\sgaStableTarget{sga6.e10.eq.5-2-3}@@M8@@
where @@M9@@ is the canonical homomorphism of 1.2.3. This follows immediately from the exact sequence
@@M10@@

\item Under the hypotheses of 5.1 and 5.2, one has, for every @@M11@@,
\sgaStableTarget{sga6.e10.eq.5-2-4}@@M12@@
which justifies the notation @@M13@@ introduced below. Indeed, one has
@@M14@@
\end{enumerate}

\sgaStableTarget{sga6.e10.sec.5.3}\subsection*{5.3.}
Composing @@M15@@ with the canonical epimorphism
@@M16@@
gives the ``first Chern class'' on @@M17@@:
\sgaStableTarget{sga6.e10.eq.5-3-1}@@M18@@

